% ============================================================================
%  Chapter 1: Introduction
% ============================================================================
\section{Introduction}
\label{sec:intro}

\paragraph{Background: from analysis tool to digital-twin counterparty.}
Large language models (LLMs) and tool-use capabilities have made AI agents able to plan
autonomously, call external tools, and execute multi-step tasks. The decisive shift is that an
agent is no longer merely a system that \emph{answers questions}; it can make economically
consequential commitments on behalf of its owner --- negotiating prices, signing terms, initiating
transfers, and settling smart contracts. When an agent holds funds and can complete payments, it
effectively constitutes a ``digital twin'' trading counterparty of its human owner.

This trend has given rise to agentic commerce and agent-to-agent (A2A) finance as research
categories \cite{a2a-finance, agent-rep}. Prior work observes that AI agents are transitioning from
analysis tools into counterparties, generating a demand for financial-grade infrastructure for
identity, authorization, payment, verification, reputation, and accountability. However,
widespread adoption amplifies three systemic risks: \emph{identity fraud} --- anyone can cheaply
mass-produce ``fake agents''; \emph{mutual-trust deficit} --- an anonymous counterparty cannot
tell whether the other side is trustworthy; and \emph{accountability vacuum} --- after an agent
misbehaves there is no clear path to attribute fault, present evidence, and impose sanctions.

\paragraph{Problem statement.}
We focus on three interconnected questions:

\begin{itemize}[leftmargin=2em,itemsep=1pt]
  \item \textbf{Q1 (Trust).} When agents transact on behalf of humans, how can mutual trust be
  established so that ``scammer agents'' do not infiltrate?
  \item \textbf{Q2 (Accountability).} When a scammer agent commits fraud, how are fault, evidence,
  and sanctions determined?
  \item \textbf{Q3 (Incentives).} Which economic and governance mechanisms make the trust
  infrastructure self-sustaining and scalable?
\end{itemize}

These questions restate the classic trust problem in a machine-to-machine, anonymity-by-default
setting: trust must come not from knowing a counterpart's history but from a verifiable, auditable,
and traceable technical infrastructure.

\paragraph{Our approach: a Schelling-point reputation community.}
We design and implement a decentralized guarantee and arbitration layer for agent-to-agent
commerce. The system is organized along the identity--trust--accountability (ITA) architecture
(Section~\ref{sec:ita}): a registry issues ERC-721 agent identities bound to real principals; a
guaranteed escrow escrows trade funds against guarantor stakes; disputes are adjudicated by a
community of staked voters whose majority --- by a Schelling-point argument --- is incentivized to
vote truthfully; and a reputation hub records outcomes on-chain and prices future guarantees. The
central novelty is the \emph{closed loop}: verdicts update reputation, reputation prices the
guarantee, and the guarantee determines whether low-reputation agents can transact at all
(Section~\ref{sec:pricing}).

The mechanism is instantiated by four production Solidity contracts (Section~\ref{sec:contracts})
with 38 passing tests, and its game-theoretic claims are validated by Monte-Carlo simulation
(Section~\ref{sec:eval}).

\paragraph{Contributions.}
\begin{itemize}[leftmargin=2em,itemsep=1pt]
  \item \textbf{A closed-loop mechanism} connecting adjudication, guarantees, and reputation --- the
  first, to our knowledge, to tie all three into one on-chain cycle (Sections~\ref{sec:mech},
  \ref{sec:pricing}).
  \item \textbf{A formal game-theoretic analysis} of the staked-voting mechanism: honest voting is
  an incentive-compatible equilibrium under a two-thirds supermajority rule; a coordinated
  malicious minority of at most $35\%$ of voters is strategically harmless; Sybil attacks are
  unprofitable because adversarial voting power scales linearly in expenditure; and the mechanism
  conserves funds exactly (Section~\ref{sec:analysis}).
  \item \textbf{A production implementation} --- four verified-style Solidity contracts, 38 tests,
  an end-to-end demo flow, and Docker one-command bootstrapping --- together with a
  \textsf{mesa}-based simulation that empirically validates the theory (Section~\ref{sec:impl}).
  \item \textbf{An honest evaluation} of the MVP's limitations (free-abstention refunds, integer
  dust, threshold approximation) and the deployment path to a public testnet
  (Sections~\ref{sec:eval}, \ref{sec:deploy}).
\end{itemize}

\paragraph{Organization.}
Section~\ref{sec:related-work} surveys related work. Section~\ref{sec:ita} presents the ITA
architecture and its mapping to the four contracts. Section~\ref{sec:mech} formalizes the mechanism
and proves its incentive properties. Section~\ref{sec:impl} describes the implementation and its
empirical evaluation. Section~\ref{sec:conclusion} concludes.
